% !TeX root = ../papers/nc-p3.tex
\section{非横断的な印付き零点からの楕円認識}
\label{sec:marked-gram-smoothing}
階数6の隣接Gram表示$\beta$とその非退化計量を用いる。
零点対を両基底の先頭に取り、$u_j=\beta_{0j}$、$v_i=\beta_{i0}$、
$T_{ij}=\beta_{ij}$ ($1\le i,j\le3$)と書く。
ここで$\beta_{00}=0$、$U=\langle u_j\rangle$、$V=\langle v_i\rangle$である。

\begin{lemma}[単純な印の三行列系の閉包]\label{lem:independent-marked-closure}
$U,V$が相補的で、両側で指定文字の一次倍四次元上に反対切断が単射なら、
その射影Gramテンソルは楕円鉛筆の閉包に属する。
\end{lemma}
\begin{proof}
直交分解後の第\ref{sec:marked-opposite-rank}節の二つの無跡三列は、
この単射性からいずれも独立である。九つの方程式$D_B=-D_A^{\mathsf t}$に対して
どの成分も次元21以上であり、階数2の解は定理\ref{thm:simple-zero-smooth-elliptic}の証明により
次元20以下である。階数1の解の次元も小さいことを確認する。
同じ証明の中心化代数の議論により、$A$の三次元空間はある共通固有ベクトル$v$を持つ。
無跡対称三行行列に対する$Xv\in kv$は二つの独立な線型条件なので、この三次元空間は$v$で決まる。
任意の対称行列で$Xv$を任意に指定でき、スカラーを引いても$V/kv$での値は変わらないからである。
二つの行列の交換子は$v$を殺し、歪対称行列と三次元ベクトルの同定で$D_A$の像は$kv$となる。
$v$が等方な場合も同じである。固定した$v$で基底の選択は九次元で、
余因子による列変換は像直線が固定された全ての非零な階数1行列（三次元）に推移的に作用する。
従って指定した階数1の値に対する三列のファイバーは六次元である。
階数1の値は高々五次元なので、対$(A,B)$は高々17次元である。
階数0は独立な可換三列を要し不可能である。
従って各成分の一般点で$D_A$は可逆で、全解は命題\ref{prop:marked-net-parameters}の
可逆チャートの閉包に入る。このチャートを基底付きで飽和した閉包と楕円鉛筆の閉包の一致は
定理\ref{thm:simple-zero-smooth-elliptic}の証明の最後による。
\end{proof}

$0\le i<k\le3,\ 0\le j<l\le3$のGram方程式を
$G_{ik,jl}=\langle\beta_{ij},\beta_{kl}\rangle-\langle\beta_{il},\beta_{kj}\rangle$とする。
$i=j=0$の九式は$U\perp V$を表す。他の27式を$T$の54座標で微分した行列を$J_T$とする。

\begin{theorem}[非横断点を通る相対Jacobi判定]\label{thm:marked-gram-smoothing}
隣接階数6、反対階数8、両側の反対核が滑らかな四次完全交叉を定めるとする。
ある零点対で$\operatorname{rank}J_T=27$なら、テンソルは基点自由楕円鉛筆の表示を持ち、
全次数で標準Koszul AS正則である。特に$R=U\cap V$の中に
\[
 C_r=(\langle r,T_{ij}\rangle)_{1\le i,j\le3},\qquad \det C_r\ne0
\]
を満たす$r$があれば、この判定が通る。
\end{theorem}
\begin{proof}
命題\ref{prop:smooth-apolar-defect}から$U,V$は三平面で$V=U^\perp$である。
$U$をGrassmann多様体上で動かし、$U,U^\perp$に基底を選ぶ空間を台とする。
これは滑らかで既約であり、非退化三平面が稠密である。
台の上で$T$を自由に動かす54次元アフィン束に残りの27式を課す。
指定の階数なら相対Jacobi判定により解空間から台への射はその点で滑らかである。
ゆえにこの解は$U$が非退化な解の閉包にある。
両側とも指定文字の一次倍へ制限した反対切断が単射となる開条件はもとの点で成立する
（滑らかな反対核には可約二次式がない）。同じ開部分内で近傍を取れる。
補題\ref{lem:independent-marked-closure}により楕円閉包への所属が特殊点にも残り、
系\ref{cor:smooth-apolar-specialization}から目的の表示が従う。

最後の十分条件を示す。27式のうち一方の指定文字だけを含む18式は$T$で線型である。
それらの微分が独立なことを確かめる。行関係の第一組は三階テンソルの係数として、
第一添字が$i$、第二添字が$j$、最終因子が$H$の元となり、
最終因子は$U$内で、$u_j$による同定で第二・第三添字について交代的である。
もう一組では最終因子が$V$内で、$v_i$による同定で第一・第三添字について交代的である。
二つの像の共通元は最終因子が$R$に入り、交代性から第二・第一因子も
各々対応する$R$の逆像に入る。従って三因子を$R$で同定した共通元は完全交代的となる。
$\dim R\le2$なので零である。各組だけの行は基底の独立性により独立である。

$\dot T_{ij}=h_{ij}r$とすればこの18式の微分は消える。残り九式への微分は、
二次小行列式の写像を$C_r$で微分して行列$(h_{ij})$に適用したものである。
三行可逆行列の余因子写像の微分は可逆である。恒等行列では
$h\mapsto\operatorname{tr}(h)I-h^{\mathsf t}$であり、他の点には基底変換で移せる。
従って残る九式にも全射である。
\end{proof}

\path{computations/verify_smooth_apolar_defect.py}は$y^2=x^3-x$上で無限遠点を$O$、
$P=(0,0)$として$L=\mathcal O(4O), M=\mathcal O(3O+P)$を取る。
鉛筆を$s$と$x^2y+y/x+11$で張り、印は$(s,1)$とする。
$s=1,x-2,y+1$で接空間次元は順に2,1,0となり、全てで$J_T$の階数は27である。
中央例の根基直線では$C_r$の階数が2なので、27式の検査は可逆な$C_r$を探すだけよりも広い。

\begin{corollary}[接空間二次元の印を持つ全候補]\label{cor:defect-two-decision}
同じ滑らかな反対核の仮定で、接空間次元2の零点対を一つ持つとする。
この印の根基$R$からの三行行列の鉛筆に階数1の元があれば標準ASではなく、
それ以外では楕円鉛筆表示が存在して標準Koszul AS正則である。
接空間次元1の印でも、根基直線の行列が階数1なら同じ排除が成立する。
\end{corollary}
\begin{proof}
$r\ne0$は$H$におけるGram像全体を消さない。対応する六式の線型結合は、
印付きの第一行と第一列が零で残る三行三列が$C_r$である。
階数1なら純テンソル関係があり、次の反転窓ではその転置が関係になる。
定理\ref{thm:no-consecutive-pure-relations}により標準ASと両立しない。
$R$が二次元で階数1の元がない場合、可逆な元があれば前定理を使える。
残るのは全ての非零元が階数2の行列鉛筆である。

このような三行鉛筆の型を確認する。$\mathbb P^1$上の写像
$\mathcal O^3\to\mathcal O(1)^3$の核を$\mathcal O(-e)$、余核を$\mathcal O(1+f)$とすれば、
両端で階数一定なので$e,f\ge0$かつ$e+f=2$である。
$e=0$なら共通核を除くと$\mathcal O(-1)^2\hookrightarrow\mathcal O^3$で余核は$\mathcal O(2)$となる。
商を与える三二次式は核に定数切断がないので独立で、$s^2,st,t^2$の基底に揃えられる。
その二本の一次関係を取れば、鉛筆は下の左行列に同値である。$f=0$は転置である。
$e=f=1$なら右・左の核直線を示す二つの一次ベクトルの成分はそれぞれ一次式全体を張る。
両方を$(-t,s,0)$に揃えると、最初の二列と最初の二行がそれぞれ$s,t$倍となり、
第三列から他の二列の倍数を引き、定数を正規化して右の形となる。
\[
 \begin{pmatrix}s&0&0\\t&s&0\\0&t&0\end{pmatrix},\qquad
 \begin{pmatrix}0&0&s\\0&0&t\\s&t&0\end{pmatrix}.
\]
各係数行列を$C_1,C_2$とする。根基内の変分だけで得られる九式の微分は、
二次小行列式写像の$C_1,C_2$での二つの微分の像の和である。
左の形とその転置では階数9、右では階数8で、右で不足する成分は左上二行二列の小行列式である。
これは二行・二列を独立にとり
$C_{ij}h_{kl}+h_{ij}C_{kl}-C_{il}h_{kj}-h_{il}C_{kj}$の係数を並べた行列の行消去で得る
（\path{computations/verify_smooth_apolar_defect.py}）。左の形ではこれで十分である。

右の形の不足する一式を埋める。補空間の添字空間を$I,J$と書くとこの形は
$C_\alpha=a_\alpha\otimes d+c\otimes b_\alpha$ ($\alpha=1,2$)で、
$a_1,a_2,c$および$b_1,b_2,d$は各双対空間の基底である。
$R$の基底$r_1,r_2$の逆像を軸の同型$J\to U,I\to V$でそれぞれ$j_\alpha,i_\alpha$とする。
線型Gram式から
$C_1(-,j_2)=C_2(-,j_1)$、$C_1(i_2,-)=C_2(i_1,-)$となる。
独立性により$J_R=\langle j_1,j_2\rangle=\ker d$、
$I_R=\langle i_1,i_2\rangle=\ker c$となり、
$a_1(i_2)=a_2(i_1),\ b_1(j_2)=b_2(j_1)$でもある。
不足した一式は$I_R\times J_R$の二行二列に対応する。
$j_3,i_3$を補い、対応する$u_3,v_3$の自己対合を1に規格化する。
$P=\langle u_3,T\rangle, Q=\langle v_3,T\rangle$とすると線型式は
\[
 P(i,j_\alpha)=C_\alpha(i,j_3),\qquad Q(i_\alpha,j)=C_\alpha(i_3,j)
\]
を与える。特にこの小矩形で$P$は列$a_\alpha|_{I_R}d(j_3)$を持ち可逆である。
ここで$\dot C_\alpha=a_\alpha\otimes d$とする。
先の二つの対称条件を満たし、$\dot P,\dot Q$を表示式に従って決めれば
線型18式を保つ変分に持ち上がる。実際、各$\dot T_{ij}$に対して指定するのは
独立な$r_1,r_2,u_3,v_3$との対合だけで、$\dot P$の第三列・$\dot Q$の第三行は自由である。

小矩形で$C_\alpha,\dot C_\alpha$はともに零だから、$T,\dot T$の各値は
$R^\perp=U+V$に入る。この矩形のGram式は$\det P+\det Q=0$であり、
構成した変分はその上で$\dot P=P,\dot Q=0$となる。その微分は$2\det P\ne0$である。
従って27式全体で階数27となり、前定理が適用できる。
\end{proof}

\begin{lemma}[根基直線上の階数2の印の平滑化]\label{lem:defect-one-rank-two}
前定理の仮定で$\dim R=1$かつ$C_r$が階数2なら、テンソルは楕円鉛筆の閉包に入る。
\end{lemma}
\begin{proof}
$r$を根基の生成元とし、両軸の$r$に対応する添字を$i_0\in I,j_0\in J$と書く
（指定した零点の文字を除いた三次元空間内の添字である）。
$U/kr,V/kr$の直交標準基底を持ち上げ、$r$に対する双曲対のもう一方の元をこれら四元と直交させる。
$T$の九個の値の座標を、この基底との対合として$C,H,P_1,P_2,Q_1,Q_2$の六行列で表す。
線型18式を解くと、$C,H$は自由である。$J=kj_0\oplus J'$とし、
二次元の$J'$を$U/kr$の上の基底で同定すると、各$i$に対して
\[
 P(i,j_0)=C(i,-)|_{J'},\qquad P(i,j')=S_i j'\quad (j'\in J')
\]
である。右辺の共変ベクトルは直交計量でベクトルと同一視し、
$S_i$は任意の対称二行行列である。$Q$側も$i,j$を交換した同形の九パラメータである。
よって自由な座標は$9+9+9+9=36$個となる。
残り九式は$C$での余因子微分を$H$に適用したものと、
各$P_\alpha,Q_\alpha$の二次小行列式との和である。

軸を固定し$C$の階数が2以上の解空間の次元が高々27であることを示す。
可逆なら$H$が一意で、$C,P,Q$の自由度が27である。
階数2なら$C$の選択は八次元で、余因子微分は階数6・核三次元である。
実際$C=\operatorname{diag}(1,1,0)$にすれば直接分かる。
左・右の核の生成元を$i_*,j_*$とし、二つの商を$\bar I=I/ki_*,\bar J=J/kj_*$とする。
この微分の余核の三座標は、九式のうち行対$(i,i_*)$、列対$(j,j_*)$の値を
$\bar I,\bar J$上の二行二列として読み、$kC$を法として取るものである。
$H$で解けるためにはこの三座標が全て零でなければならない。
従って条件を課す前の上界$8+9+9+3=29$を少なくとも二つ下げればよい。

$j_*=j_0$（直線の一致）なら、階数2の右核が指定直線となること自体が余次元2である。
以下$j_*,j_0$は独立とする。$S_{i_*}$がスカラー行列なら九個の$P$の自由座標に
独立な二条件がかかる。残る場合は$C,Q,S_{i_*}$を固定する。
$j_*$の$J'$成分を$k$と書き、それと独立な$l\in J'$を選ぶ。
$j_0,l$の像は$\bar J$の基底である。$i_*$を補う二つの$i$について$S_i$を独立に変える。
先の二行二列のうち$j_0$の列は変わらない。
実際$P(i,j_0)$は固定され、$P(i_*,j_0)=0$である。
$l$の列の$i$行の変化は
\[
  \langle \dot S_i l,S_{i_*} k\rangle-\langle\dot S_i k,S_{i_*} l\rangle
\]
である。$k,l$は基底で、全ての対称二行行列と可換な行列はスカラーに限られるので、
これは$\dot S_i$の非零の線型形式である。二行を独立に指定でき、
$C(-,j_0)\ne0$なので得られる二次元空間は$kC$を法としても単射に写る。
固定した他の座標のもとで条件は$S_i$についてアフィン線型である。従って少なくとも二条件がかかる。
これで全ての場合に固定軸の上界27が従う。

二つの軸を動かす台は先のGrassmann多様体の基底束である。
根基直線を持つ三平面はその台で余次元1以上である（非退化条件の行列式が恒等零ではない）。
根基が二次元以上の閉部分を除き、根基が直線の閉境界上で$C$の階数が1以下となる
閉条件も除いてよい。指定点は残る。54変数に27式を課した全解空間の成分の次元は
台の次元に27を加えた値以上である。一方、先の上界により残った閉境界だけでは
この次元に達しない。従って指定点は相補三平面を持つ解の閉包にある。
両軸での単射性の開条件を保ち、補題\ref{lem:independent-marked-closure}を用いればよい。
\end{proof}

\begin{theorem}[滑らかな反対四次曲線を持つ反転四窓の分類]\label{thm:smooth-apolar-complete}
両対称テンソルの隣接切断が階数6、反対切断が階数8で、
両反対核が滑らかな四次完全交叉を定めるとする。付随する反転四窓が標準Koszul AS正則となる
必要十分条件は、第一窓の六次元関係空間が非零の純テンソルを含まないことである。
この条件は全てのAS点が基点自由楕円鉛筆から得られることとも対応し、各点で
二つの曲線の同定と順序付きの二束・鉛筆が復元される。
\end{theorem}
\begin{proof}
純テンソルのある場合はその転置と連続するから既出の成長障害を使う。
純テンソルがない場合、六つの$(1,1)$式には少なくとも一つ零点対がある。
実際、六次元Segre多様体を射影空間で六つの超平面と交わらせている。
印の接空間次元は命題\ref{prop:smooth-apolar-defect}により$0,1,2$のいずれかである。
単純点では定理\ref{thm:simple-zero-smooth-elliptic}を使う。
次元1なら根基の行列は非零で階数1ではない。階数3は定理\ref{thm:marked-gram-smoothing}、
階数2は補題\ref{lem:defect-one-rank-two}と系\ref{cor:smooth-apolar-specialization}による。
次元2は系\ref{cor:defect-two-decision}で尽くされる。いずれも滑らかな曲線と基点自由鉛筆へ戻るので、
全次数の正則性とデータの一意性は定理\ref{thm:elliptic-pencil-reconstruction}による。
\end{proof}

\begin{corollary}[反転窓の反対階数10の排除]\label{cor:bisymmetric-high-rank-exclusion}
隣接階数6の両対称Gramテンソルで、両反対核がいずれも非零の一次式の積を含まないとする。
付随する反転窓が標準ASならテンソルは楕円鉛筆の閉包にあり、反対階数は8以下である。
特に反対階数10の反転窓は全て排除される。階数9の標準AS反転窓があれば、
少なくとも一方の一次元反対核は階数2以下の二次式で生成される。
\end{corollary}
\begin{proof}
指定文字の一次倍上の反対切断の単射性と、平方が核に入らないことは
両反対核についての仮定から従い、階数8や四次曲線の滑らかさを必要としない。
従って命題\ref{prop:smooth-apolar-defect}の三平面と接空間次元の議論が通る。
標準ASなら根基の階数1の元は純テンソル障害により存在しない。
単純点は補題\ref{lem:independent-marked-closure}、それ以外の印は前定理の
二つの平滑化証明を、楕円鉛筆の閉包を得る段階まで適用する。
閉包では反対切断の九次小行列式が消える。階数10では両反対核が零、
階数9ではそれぞれ一次元なので残る二主張を得る。
\end{proof}
